\documentclass{article}
\usepackage{ctex}
\usepackage{amsmath}
\usepackage{tikz}
\usepackage{unicode-math}
\usetikzlibrary{arrows.meta, decorations.markings,patterns}
\begin{document}

2.3-4平行板电容器(极板面积为$S$,间距为$d)$中间有两层厚度各为$d_1$和$d_2(d_1+d_2=d)$,介电常量各为$\varepsilon_1$和$\varepsilon_2$的电介质层,试求：

(1)电容$C$

(2) 当金属极板上带电面密度为±$\sigma_{e_0}$时,两层介质间分界面上的极化电荷面密度 $\sigma_e'$

(3)极板间电势差$U$

(4)两层介质中的电位移$D$





\newpage

2.3-7 一平行板电容器两极板相距为$d$,面积为$S$,电势差为$U$,其中放有一层厚为$t$ 的介质,介电常量为 $\varepsilon$,介质两边都是空气. 略去边缘效应.求:

(1)介质中的电场强度$E$,电位移$D$和电极化强度$P$

(2)极板上的电荷量$Q$

(3)极板和介质间隙中的电场强度$E$

(4)电容$C$









\newpage

2.3-12 一平行板电容器两极板相距为$d$,其间充满了两部分介质(左右放置),左侧介电常量为$\varepsilon_1$ 的介质所占的面积为 $S_1$,右侧介电常量为$\varepsilon_2$的介质所占的面积为$S_2$略去边缘效应,求电容$C$
















\newpage

2.3-13 一平行板电容器两极板的面积都是 S,相距为 $d$,今在其间平行地插入厚度为 $t$、介电常量为$\varepsilon$的均匀介质 ,其面积为 $\frac{S}{2}$. 设两板分别带电荷$Q$和-$Q$,略去边缘效应,求：

(1)两板电势差 $U$

(2)电容$C$

(3)介质的极化电荷面密度 $\sigma_e'$



















\newpage

2.3-19 一半径为$R$ 的导体球带电荷 $Q$,球外有一层同心球壳的均匀电介质,其内外半径分别为$a$ 和 $b$,介电常量为$\varepsilon$.求：

(1)介质内外的电场强度$E$ 和电位移 $D$

(2)介质内的电极化强度$P$和表面上的极化电荷面密度$\sigma_e'$

(3)介质内的极化电荷体密度$\rho_e'$ 为 多 少 ? 















\newpage

2.3-24 圆柱形电容器是由半径为$a$ 的导线和与它同轴的导体圆简构成,圆筒内半径为$b$,长为$l$,其间充满了两层同轴圆筒形的均匀介质,
分界面的半径为$r$,介电常量分别为$\varepsilon_1$和$\varepsilon_2$,略去边缘效应,求电容$C.$













\newpage

2.3-35 两共轴的导体圆筒,内简外半径为$R_1$,外简内半径为$R_2(R_2<2R_1)$,其间有两层均匀介质,分界面的半径为$r$,
内层介电常量为 $\varepsilon _1$,外层介电常量为 $\varepsilon_2=\varepsilon_1/2$,两介质的介电强度都是 $E_m.$ 

当电压升高时,哪层介质先击穿？

证明：两筒最大的电势差$U_m=\frac{E_m r}{2}\ln\frac{R_{2}^{2}}{rR_{1}}$










\newpage

2.3-39 一平行板电容器极板面积为 $S$，间距为 $d$，接在电源上以维持其电压为 $U$．将一块厚度为 $d$、介电常量为 $\varepsilon$ 的均匀电介质板插入极板间空隙．计算：

(1) 静电能的改变 ; 

(2)电场对电源所做的功；

(3) 电场对介质板做的功.













\end{document}